\chapter{样式速查（编写参考）}
\label{ch:gallery}

本附录汇总全书可用排版单元，编写时查阅。编号全自动，引用一律 \verb|\cref{...}|。

\section{公式与编号}
公式按「章.序」编号（已设 \verb|\numberwithin{equation}{chapter}|）：
\begin{equation}\label{eq:demo}
  \mathbf{R}\mathbf{R}^\top=\mathbf{I},\qquad \det\mathbf{R}=1 .
\end{equation}
引用示例：正交性见 \cref{eq:demo}；不编号用 \verb|equation*| 或 \verb|\notag|。

\section{脚注与行内批注}
脚注示例\footnote{footmisc 脚注，置于页面底部。}。行内批注（小字彩色、不入页边——旁注会溢出版心，已弃用）：\pz{蓝色批注}、\pzr{红色重点}、\rebuilt{重建待核对}。

\section{内置彩色盒（原生）}
\begin{definition}[特殊正交群]\label{def:demo}
  绿色 \texttt{definition}（main 色）。
\end{definition}
\begin{theorem}[指数映射]\label{thm:demo}
  橙色 \texttt{theorem}（second 色）。引用：\cref{def:demo}、\cref{thm:demo}。
\end{theorem}
\begin{proposition}
  蓝色 \texttt{proposition}（third 色）。
\end{proposition}
\begin{note}
  \texttt{note} 提示盒。
\end{note}
\begin{example}
  \texttt{example} 例题盒（独立计数）。
\end{example}

\section{自定义单元（\texttt{\textbackslash elegantnewtheorem}）}
\begin{algo}[ESKF 预测]\label{alg:demo}
  蓝色「算法」盒，定义见 \verb|styles.tex|。引用：\cref{alg:demo}。
\end{algo}
\begin{insight}
  绿色「洞见」盒，用于强调关键直觉。
\end{insight}

\section{代码框与推导框（区分内容类型）}
代码用 \texttt{codebox}：
\begin{codebox}[C++]{Sophus 构造 SO(3)}
Eigen::Matrix3d R = Eigen::AngleAxisd(M_PI/4,
                      Eigen::Vector3d(0,1,0)).toRotationMatrix();
Sophus::SO3d SO3_R(R);   // 从旋转矩阵构造
\end{codebox}

推导用 \texttt{derivation}（可跨页）：
\begin{derivation}[罗德里格斯公式]
由 $\exp(\theta\mathbf{a}^\wedge)=\sum_{n\ge 0}\tfrac{1}{n!}(\theta\mathbf{a}^\wedge)^n$，
利用 $\mathbf{a}^\wedge\mathbf{a}^\wedge=\mathbf{a}\mathbf{a}^\top-\mathbf{I}$ 与 $(\mathbf{a}^\wedge)^3=-\mathbf{a}^\wedge$，整理得
\begin{equation*}
  \exp(\theta\mathbf{a}^\wedge)=\cos\theta\,\mathbf{I}+(1-\cos\theta)\,\mathbf{a}\mathbf{a}^\top+\sin\theta\,\mathbf{a}^\wedge .
\end{equation*}
\end{derivation}

\section{内置单元清单}
彩色定理类：\texttt{theorem / definition / lemma / corollary / proposition / axiom / postulate}；
练习类：\texttt{example / exercise / problem}；
说明类：\texttt{proof / solution / note / remark / assumption / conclusion / property / custom\{标题\}}；
版块类：\texttt{introduction / problemset}。星号版（\texttt{theorem*} 等）不编号。
